\begin{abstract}
Enterprise Model Context Protocol (MCP) agents frequently expose large, deeply constrained JSON Schema specifications that inflate prompt token budgets. While lazy schema resolution seeks to mitigate context bloat, Learned Adaptive Schema Resolution (LASR-MCP) operates by pruning upfront schemas and dynamically injecting violated constraints upon validation failure. However, preliminary v1.0 evaluations suffered from heuristic token double-counting and neglected the multi-turn cost of retry cascades. In this paper, we present a rigorous empirical and theoretical evaluation of LASR-MCP (\texttt{mcp-inventory-bench} v2.0.3) across $N=3,840$ completed trials across 32 tasks, six experimental arms, and 20 fixed-seed executions per task-arm pair ($7,860$ response-bearing API invocations on \texttt{gemini-3.5-flash-lite} at temperature $0.2$, alongside 154 client transport retries). Using provider-reported prompt-plus-candidate ($P+C$) token usage from \texttt{usageMetadata} and a paired task-cluster percentile bootstrap (10,000 resamples, seed 42), we establish three foundational findings. First, naive zero-constraint lean schemas (\texttt{ci\_k0}) induce severe multi-attempt retry cascades ($0.0\%$ first-pass success), degrading final structural reliability to $87.5\%$ and inflating net token consumption by $84.3\%$ ($2,057.7$ vs.\ $1,116.8$ mean recorded $P+C$ tokens for Full Control; $95\%$ signed-savings CI: $[-125.15\%, -49.21\%]$). Second, prioritized upfront bundling (\texttt{ci\_k2\_reactive}) slashes net token consumption by $31.7\%$ compared to \texttt{ci\_k0} ($1,406.5$ mean recorded $P+C$ tokens), recovers $96.9\%$ final structural reliability ($93.8\%$ semantic reliability), and operates within $2.1\%$ of a task-aware oracle reference baseline (\texttt{oracle\_reactive}, $1,377.8$ tokens, where LASR $k=2$ is $2.08\%$ higher using the oracle as denominator). Third, we characterize an observed \emph{Retry-Amortization Boundary} across the four evaluated enterprise task archetypes: across the full benchmark, Full Control achieves a mean recorded $P+C$ subtotal of $1,116.8$ tokens per trial versus $1,406.5$ for LASR $k=2$, meaning LASR $k=2$'s overall workload is approximately $25.94\%$ higher than Full Control's ($95\%$ CI on workload increase: $[+3.16\%, +52.55\%]$, corresponding to signed relative savings of $-25.94\%$, $95\%$ CI: $[-52.55\%, -3.16\%]$, where the 95\% bootstrap CI excludes zero); we do not describe the overall benchmark result as token savings. Rather, token savings are strictly cluster-specific: in routine operational regimes ($C_1$) where active task constraints fall within the upfront budget ($m_{\text{task}}(x) \le k$), LASR $k=2$ achieves $75.0\%$ first-pass success and delivers approximately $13.5\%$ lower $P+C$ usage ($847.6$ vs.\ $979.8$ tokens) at $100.0\%$ reliability; conversely, in deep inter-field constraint regimes ($C_2$--$C_4$, $m_{\text{task}}(x) > k$), single-tool isolation favors monolithic schemas by avoiding multi-turn prompt re-transmission. We distinguish the full multi-session collection timeline from the final resumed execution session and acknowledge that all 20 repetitions per task-arm produced identical outputs under the repeated fixed seed. All raw telemetry, audit scripts, and cryptographic checksums are publicly released to establish an auditable benchmark for agentic tool use.
\end{abstract}
